\documentclass[12pt]{article}
\usepackage[margin=0.6in,top=1.15in,bottom=0.55in]{geometry}
\usepackage{lmodern}
\usepackage{microtype}
\usepackage{titlesec}
\usepackage{enumitem}
\usepackage[hidelinks]{hyperref}
\usepackage{../../latex/grader-worksheet}

\setlength{\parindent}{0pt}
\setlength{\parskip}{0.25em}
\titleformat{\section}{\large\bfseries\scshape}{}{0pt}{}

\GraderSetup{assignment-id=U1L2LE,template-revision=1}

\begin{document}

\begin{center}
{\LARGE\bfseries Week 2 Lit Example Worksheet}\par\vspace{0.05em}
{\large\scshape Macbeth: Shifting Words, Shifting Reasons}\par\vspace{0.12em}
\rule{0.78\textwidth}{0.5pt}
\end{center}

\textbf{Purpose:} Apply Week 2's Whately method to Macbeth. Define important terms, separate different senses of the same word, and test whether the reasoning still works.

\textbf{Directions:} Use the Lit Example Reader. Quote only briefly. Write in complete, precise sentences. Do not merely say that Shakespeare is using a theme or symbol; explain how a word's meaning affects an inference.

\section*{Part 1: Fair and Foul}

\textbf{Question 1.1}\\[-0.15em]
In the first passage, identify two possible meanings of \textit{fair} and two possible meanings of \textit{foul}. Then explain why confusing moral meaning with outward appearance or favorable outcome could mislead Macbeth.

\GraderAnswerBox{Part 1 Q1}{2.25in}

\newpage

\textbf{Question 1.2}\\[-0.15em]
Macbeth says, ``So foul and fair a day I have not seen.'' Does this line by itself prove a contradiction, or can the words be using different senses? Explain the logical difference.

\GraderAnswerBox{Part 1 Q2}{2.15in}

\newpage

\section*{Part 2: What Counts as a Man?}

\textbf{Question 2.1}\\[-0.15em]
Reconstruct Macbeth's use of \textit{man}: what does he mean when he says, ``I dare do all that may become a man; / Who dares do more is none''? State the term, Macbeth's sense of it, and the limit he is drawing.

\GraderAnswerBox{Part 2 Q1}{2.35in}

\textbf{Question 2.2}\\[-0.15em]
Reconstruct Lady Macbeth's use of \textit{man}: what does she imply when she says Macbeth was a man when he dared the murder and would be ``so much more the man'' if he did it? State her sense of the term and explain whether her argument depends on a shift in meaning.

\GraderAnswerBox{Part 2 Q2}{2.45in}

\newpage

\section*{Part 3: Done and Settled}

\textbf{Question 3.1}\\[-0.15em]
Macbeth says, ``If it were done when 'tis done.'' Identify at least two senses of \textit{done} in this passage. Then explain why physically completing an action does not prove that its consequences or moral judgment are finished.

\GraderAnswerBox{Part 3 Q1}{2.65in}

\section*{Part 4: Testing an Equivocation}

\textbf{Question 4.1}\\[-0.15em]
Choose either \textit{fair/foul}, \textit{man}, or \textit{done}. Build a short argument that would become unreliable if that word changed meaning halfway through. Then rewrite the argument with the term clarified so the shift is visible.

\GraderAnswerBox{Part 4 Q1}{2.75in}

\end{document}
